\documentclass{umk-matematika}

\begin{document}

\maketitlepage
    {Практическая работа №28}
    {Конус. Площадь поверхности и объём}
    {}

\section{Фундамент (1--4 балла)}

\textbf{Задача 1 (1 балл).} Радиус основания конуса $6$ см, образующая $10$ см. Найдите площадь боковой поверхности.

\textbf{Задача 2 (2 балла).} Радиус основания конуса $4$ см, высота $3$ см. Найдите объём.

\textbf{Задача 3 (3 балла).} Площадь боковой поверхности конуса $45\pi$ см², радиус основания $5$ см. Найдите образующую.

\textbf{Задача 4 (4 балла).} Объём конуса $100\pi$ см³, радиус основания $5$ см. Найдите высоту.

\section{Стены (5--7 баллов)}

\textbf{Задача 5 (5 баллов).} Осевое сечение конуса — равносторонний треугольник со стороной $8$ см. Найдите площадь полной поверхности конуса.

\textbf{Задача 6 (6 баллов).} Образующая конуса $13$ см, высота $12$ см. Найдите объём конуса.

\textbf{Задача 7 (7 баллов).} Радиус основания конуса $7$ см, высота $24$ см. Найдите площадь полной поверхности.

\section{Кровля (8--10 баллов)}

\textbf{Задача 8 (8 баллов).} Куча песка имеет форму конуса. Длина окружности основания $6\pi$ м, образующая $5$ м. Найдите объём песка.

\textbf{Задача 9 (9 баллов).} Крыша башни — конус с диаметром основания $10$ м и высотой $12$ м. Сколько листов кровли нужно, если один лист покрывает $2$ м²? (С запасом $10\%$)

\textbf{Задача 10 (10 баллов).} Резервуар состоит из цилиндра (радиус $3$ м, высота $4$ м) и конуса (такой же радиус, высота $2$ м), поставленного сверху. Найдите общий объём.

\newpage
\section*{Ответы}

\begin{enumerate}
  \item $60\pi$ см²
  \item $16\pi$ см³
  \item $9$ см
  \item $12$ см
  \item $48\pi$ см²
  \item $100\pi$ см³
  \item $224\pi$ см²
  \item $12\pi$ м³ $\approx 37{,}7$ м³
  \item $113$ листов
  \item $42\pi$ м³ $\approx 131{,}9$ м³
\end{enumerate}

\end{document}